\clearpage
\SourceLeaf{162}
\begin{center}
[VI.]

{\large [EINE AUSGLEICHSFORMEL FÜR\\ MORTALITÄTSTAFELN.]}
\end{center}



Das reine Resultat der Erfahrungen über die Tontinen ist folgendes:

$a$ Anzahl der zu Anfang des $3^{\text{ten}}$, $\dfrac{a}{x}$ Anzahl der zu Anfang des
$n^{\text{ten}}$ Jahres Lebenden:

\begin{center}
\begin{tabular}{cc|cc}
$n$ & $\log x$ & $n$ & $\log x$ \\
3 & 0 & 52 & 0,24220 \\
7 & 0,03547 & 57 & 0,28735 \\
12 & 0,06075 & 62 & 0,34577 \\
17 & 0,07687 & 67 & 0,42323 \\
22 & 0,09718 & 72 & 0,54021 \\
27 & 0,11871 & 77 & 0,72623 \\
32 & 0,14038 & 82 & 1,02073 \\
37 & 0,16404 & 87 & 1,43402 \\
42 & 0,18579 & 92 & 2,09812 \\
47 & 0,21004 & 97 & 3,17730 \\
\end{tabular}
\end{center}

Befriedigend wird $\log x$ dargestellt durch die Formel
\[
\log x = A+Bb^n-Cc^n,
\]
